笛卡尔树同时满足：中序遍历顺序就是原数组顺序，父亲的值不大于儿子的值。
因此点 $i$ 对应数组位置 $i$，不需要重新编号。

\texttt{cartesian\_tree(a)} 返回根和儿子数组。
返回值的 \texttt{first} 是根，\texttt{second[u]} 是点 $u$ 的左右儿子；不存在的儿子为 $-1$。

建小根树时使用 \texttt{a[st.back()] > a[i]}；只把这里的 \texttt{>} 改成 \texttt{<}，就得到大根树。
当前相等元素不弹栈，所以相等时靠左的位置更高；若希望靠右的位置更高，小根使用 \texttt{>=}，大根使用 \texttt{<=}。

常用结论：区间 $[l,r]$ 的最小值位置是笛卡尔树中 $l,r$ 的 LCA。
